\section{Methods}

Information extraction typically consists of two stages: identifying
candidate phrases and determining their meaning through
classification. In this study, we focus on the phrase classification
task and study how different adaptations of LLMs cope with it.

\subsection{Model Benchmarking and Data Scaling}

Fair evaluation of LLM-based methods in this context is challenging.
Meaningful and reproducible evaluation requires a large, consistently
annotated public dataset. Because such medical datasets are rarely
available, we first evaluate model families on a well-established
non-medical benchmark dataset
MultiNERD~\cite{tedeschi-navigli-2022-multinerd} to understand
general properties and trends. This dataset contains named entity
annotations for Wikipedia and Wikinews sentences. Its diverse
annotation scheme includes 16 entity classes across more than 160,000
sentences, providing substantial variability for model evaluation.
Our main interest is how performance scales with the amount of
annotated training data and which adaptation method performs best
when ample training data are available. We report performance using
the F1 score, which balances precision and recall.

\subsection{Creation of the Target Clinical Dataset}

Following \v{S}uvalov \emph{et al.}, we consider only patient
anamnesis notes. We use MedGemma with system prompts adapted from
\cite{suvalov-2026-extract} to extract candidate drug discontinuation
events for manual review. Although originally designed for a narrower
extraction task, these prompts achieved over 90\% recall, making them
suitable for candidate generation.

The event cohort was assembled from Est-Health 30, which contains
centrally collected electronic health records from 2012--2024 for a
random 30\% sample of the Estonian
population~\cite{reisberg2026eh30}. To capture rare drug
discontinuation events, we selected records based on drug purchases
and grouped them into four subcohorts: potentially incompatible drug
combinations, prolonged treatment gaps ($\geq$ 6 months), drug
discontinuations and switches. For each cohort, we collected all
anamnesis notes recorded within six months of the index event. After
phrase extraction with MedGemma (27B), we randomly sampled 1,000
potential discontinuation phrases from each cohort, yielding a
dataset of 4,000 notes. Three experts annotated 3,300 candidate
phrases. Each anamnesis note was annotated with one or more event
labels. The labels \emph{adverse reaction}, \emph{treatment success},
\emph{treatment inefficacy}, \emph{contraindication or unsafe co-use}
and \emph{other} represent reasons for drug discontinuation, whereas
\emph{adverse reaction without discontinuation}, \emph{dosage change}
and \emph{other non-discontinuation} describe medication-related
phrases without discontinuation. Unsuitable phrases were labelled
\emph{incorrect phrase}.

\subsection{Experimental Setup and Performance Metrics}

For the MultiNERD benchmark, we evaluated model performance using
training sets of 10, 100, 1,000 and 20,000 labelled examples. For the
clinical dataset, we repeated the analysis with training-set sizes
ranging from 100 to 2,500 annotations. In all experiments, the test
set remained fixed. For the clinical task, we also collapsed the nine
event labels into a binary outcome indicating a drug discontinuation event.

The main experimental factors were model family and size, learning
method, training-set size, and prompt initialisation. For prompt
tuning, we compared how random initialisation compares with a
user-selected starting prompt describing the task. We recorded test
F1, predictions, model checkpoints, and learned prompt embeddings.

Experiments were run on a single GPU using BF16 precision and 8-bit
AdamW optimiser. We used a batch size of 8 and trained models with
early stopping based on validation F1, using learning rates of $5
\times 10^{-5}$ for fine-tuning and $10^{-2}$/$10^{-3}$ for prompt
tuning of decoder/encoder models.

\subsection{Context-Aware Phrase Classification}

The objective of context-aware phrase classification is to assign a
binary or multiclass label to a candidate phrase using both the
surrounding text and task instructions (\emph{prompt}). Annotated
examples are used to further adapt the model to the target
classification task. For autoregressive LLMs, context-aware
classification can be framed as a question-answering task in which
the model's output represents the target class. For encoder models
such as \textsc{BERT}, the task must be reduced to standard text
classification. The prompt and context can be encoded as part of the
input sequence using \textsc{[sep]} tokens for explicit structuring.
The output is passed to a linear classification head.

Prompting itself can be done by asking questions in natural language.
Supervised learning allows us to move beyond manual prompt
engineering. Parameter-efficient fine-tuning (PEFT) updates only a
small subset of model parameters. For example prompt tuning learns a
virtual (\emph{soft}) prompt prepended to the input, leaving the base
model unchanged~\cite{lester2021power}.
